
% ============================================================
% SECTION 07
% FORMS AND DOMAINS OF INVARIANCE
% Chapter 02: Invariance
% Book I: Ontology of Viability
% ============================================================

\section{Forms and Domains of Invariance}
\label{sec:forms-and-domains-of-invariance}

% ------------------------------------------------------------
% 7.1. THE NEED TO DISTINGUISH FORMS OF INVARIANCE
% ------------------------------------------------------------

The preceding sections have established a
provisional relational definition of invariance.
Invariance concerns the relation between an
existent, or a determinate aspect of an existent,
and a specified domain of changes or
transformations, within which that existent
or aspect remains unchanged in a particular
respect.

This definition identifies the general structure
of an invariance claim. It does not, however,
imply that every instance of invariance has
the same form or operates within the same
domain.

Different determinations may remain unchanged
under different kinds of transformations.
A relation may remain invariant while a
configuration changes. A quantitative
determination may remain constant while
other determinations vary. A particular
organization may persist under transformations
that alter its elements or their positions.

These differences make it necessary to
distinguish forms and domains of invariance.

The purpose of this section is not to establish
a final and exhaustive classification. Rather,
it is to identify several analytical distinctions
that help specify what is being described
when invariance is attributed to an existent
or one of its determinations.

The classification developed here remains
provisional and does not imply that the
distinguished forms constitute independent
ontological categories.

% ------------------------------------------------------------
% 7.2. THE DOMAIN OF TRANSFORMATION
% ------------------------------------------------------------

The domain of transformation specifies the
changes or transformations in relation to
which invariance is being examined.

Without such a domain, the statement that
something is invariant remains incomplete.

A determination may remain unchanged under
one set of transformations while changing
under another. Therefore, invariance cannot
be attributed without qualification.

For analytical purposes, a domain of
transformation may be specified according
to the kind of alteration being considered,
the conditions under which the alteration
occurs, or the range of transformations
included in the examination.

These specifications are methodological
requirements. They do not imply that
the domain of transformations exists as
an independent ontological entity.

The distinction can be represented
schematically as follows:

\[
\mathcal{D}_1 \neq \mathcal{D}_2
\]

where \(\mathcal{D}_1\) and \(\mathcal{D}_2\)
represent different domains of transformation.

If a determination remains unchanged
within \(\mathcal{D}_1\), this does not
establish that it remains unchanged
within \(\mathcal{D}_2\).

Accordingly, every meaningful claim
of invariance must make its domain
explicit or establish it through
the context of the analysis.

% ------------------------------------------------------------
% 7.3. INVARIANCE OF QUANTITATIVE
% DETERMINATIONS
% ------------------------------------------------------------

One possible form of invariance concerns
the persistence of a quantitative
determination under specified transformations.

A quantity may remain unchanged while
other aspects of the existent undergo
alteration.

For example, a particular quantity
may remain constant under transformations
that modify the configuration, position,
or relations of the elements involved.

This does not establish that the entire
existent remains unchanged. It identifies
only the invariance of the specified
quantitative determination.

The relevant claim can be expressed
schematically as:

\[
Q(T(A)) = Q(A)
\]

where \(Q\) represents a specified
quantitative determination and \(T\)
a transformation within the relevant
domain.

The expression is meaningful only when
the quantity and the transformation
are adequately specified.

It does not establish that quantitative
invariance is more fundamental than
other forms of invariance, nor that
every quantity can remain invariant
under transformation.

The ontological significance of
quantitative determinations remains
dependent upon the broader account
of the existent in which they occur.

% ------------------------------------------------------------
% 7.4. INVARIANCE OF QUALITATIVE
% DETERMINATIONS
% ------------------------------------------------------------

Invariance may also concern qualitative
determinations.

A qualitative determination refers here
to a specified character or feature
of an existent that is not being
considered primarily as a numerical
quantity.

Such a determination may remain
unchanged under some transformations
while other qualitative or quantitative
aspects are modified.

However, the identification of
qualitative invariance requires
careful specification.

A qualitative description may group
together different determinations
under a common term. The persistence
of that description does not
necessarily establish that the
underlying determination has remained
unchanged in every relevant respect.

The analysis must therefore distinguish
between the invariance of a determination
and the continued applicability of
a particular description.

The latter may reflect the limits
or conventions of the descriptive
framework rather than the complete
ontological status of the existent.

This distinction is particularly
important when qualitative descriptions
are used to compare different
conditions or phases.

% ------------------------------------------------------------
% 7.5. RELATIONAL INVARIANCE
% ------------------------------------------------------------

A further form of invariance concerns
relations.

A relation between determinations
or entities may remain unchanged
under transformations that alter
other aspects of the configuration.

For example, the relative ordering
of elements may remain unchanged
even when their positions or
distances are modified.

In such a case, the invariance
concerns the specified relation,
not necessarily the elements,
their properties, or the complete
configuration in which they occur.

Relational invariance must therefore
be distinguished from the invariance
of the relata and from the invariance
of the total relational organization.

The following expression represents
a relation \(R\) that remains
unchanged under a specified
transformation:

\[
R(T(A)) = R(A)
\]

This expression is schematic.
It does not determine the ontological
status of the relation or establish
whether the relation is independent
of the entities involved.

The purpose is only to represent
the possibility that a relation
may be invariant while other
determinations undergo change.

% ------------------------------------------------------------
% 7.6. STRUCTURAL INVARIANCE
% ------------------------------------------------------------

Structural invariance concerns
the persistence of a specified
organization or structural
determination under a defined
domain of transformations.

As discussed in the preceding
section, the invariance of one
structural relation does not
necessarily establish the
invariance of the entire structure.

Likewise, a structure may preserve
a specified organization even
when some of its constituent
determinations are altered.

Structural invariance therefore
requires an explicit distinction
between the organization considered
as a whole and the particular
features or relations through
which that organization is described.

This form of invariance may be
relevant to the examination of
continuity and identity, but
it cannot independently establish
either.

The classification of structural
invariance as a distinct analytical
form does not imply that structure
is an ontologically fundamental
category.

Its deeper status remains open
to the broader development of
the ontology.

% ------------------------------------------------------------
% 7.7. TEMPORAL AND NON-TEMPORAL
% DOMAINS OF INVARIANCE
% ------------------------------------------------------------

The domain of invariance may also
be considered in relation to
the presence or absence of
temporal transformation.

In a temporal domain, a determination
is examined across different
times or successive conditions.

A particular feature may remain
unchanged during a specified
interval, while changing before
or after that interval.

Such invariance is necessarily
qualified by the temporal range
under consideration.

However, not every transformation
must be understood as a temporal
change.

A comparison may concern different
configurations, perspectives,
relations, or conditions without
the analysis being primarily
about temporal succession.

For this reason, the domain
of invariance should not be
reduced in advance to time.

This distinction does not establish
a complete theory of non-temporal
invariance. It only prevents
the assumption that every
invariance claim must be
formulated as persistence
through time.

The ontological relationship
between time, transformation,
and invariance remains open.

% ------------------------------------------------------------
% 7.8. LOCAL AND GENERAL INVARIANCE
% ------------------------------------------------------------

A further distinction concerns
the scope of the determination
being examined.

Local invariance refers to the
persistence of a specified
determination within a limited
aspect, region, relation, or
domain of an existent.

General invariance refers to
a broader claim concerning
a more extensive organization
or range of determinations.

These terms describe differences
in analytical scope rather than
independent ontological categories.

A local invariant does not
necessarily establish a general
invariant.

Likewise, a claim of general
invariance requires more than
the identification of one or
several local invariants.

The distinction can be represented
schematically:

\[
I_{\mathrm{local}}
\not\Rightarrow
I_{\mathrm{general}}
\]

The expression indicates that
local invariance alone does not
logically establish general
invariance.

The same requirement applies
to claims concerning an entire
existent. The invariance of
a limited number of determinations
must not be generalized to
the whole without sufficient
grounds.

% ------------------------------------------------------------
% 7.9. INTERDEPENDENCE AND OVERLAPPING
% DOMAINS
% ------------------------------------------------------------

The forms of invariance distinguished
in this section should not be
understood as mutually exclusive.

A particular determination may
be quantitative and structural,
relational and temporal, or
qualitative and local.

The distinctions identify
different analytical aspects
of invariance rather than
separate classes that must
remain isolated.

Moreover, different domains
of transformation may overlap.

A transformation may alter
several determinations at once,
while preserving others.
Consequently, the identification
of an invariant feature may
require attention to more
than one domain or analytical
dimension.

The classification must therefore
remain flexible enough to
accommodate cases in which
different forms of invariance
are simultaneously relevant.

This flexibility does not
eliminate the need for precision.
On the contrary, the more
domains and determinations
are involved, the more important
it becomes to specify the
respect in which invariance
is being asserted.

% ------------------------------------------------------------
% 7.10. SYNTHESIS AND TRANSITION
% ------------------------------------------------------------

This section has distinguished
several analytical forms and
domains of invariance.

These include:

\begin{enumerate}
    \item Quantitative invariance.
    \item Qualitative invariance.
    \item Relational invariance.
    \item Structural invariance.
    \item Invariance within temporal
    and non-temporal domains.
    \item Local and general invariance.
\end{enumerate}

These distinctions do not
constitute an exhaustive
classification, nor do they
establish separate ontological
kinds of invariance.

They clarify that every claim
of invariance must specify
the determination under
consideration, the domain
of transformation, and the
respect in which something
remains unchanged.

The analysis has also established
that different forms of invariance
may overlap and that the
invariance of a limited aspect
cannot automatically be
generalized to the whole existent.

The next section examines
invariance and knowledge.
It will address the distinction
between the ontological status
of what remains unchanged
and the epistemic conditions
under which such invariance
is identified, described,
and justified.

% ============================================================
% END OF SECTION 07
% ============================================================
